% Plantilla base para el producto integrador.
\documentclass[12pt,letterpaper,oneside]{article}
\usepackage[margin=2.54cm]{geometry}
\usepackage[spanish,es-tabla]{babel}
\usepackage[utf8]{inputenc}
\usepackage[T1]{fontenc}
\usepackage[scaled=.96]{helvet}
\renewcommand{\familydefault}{\sfdefault}
\usepackage{setspace}
\usepackage{graphicx}
\usepackage[table]{xcolor}
\usepackage[authoryear,round]{natbib}
\usepackage{xurl}
\usepackage[hidelinks]{hyperref}

\setstretch{1.5}
\setlength{\parindent}{1.25cm}
\setlength{\parskip}{0.25em}

\begin{document}
\begin{center}
{\bfseries Universidad Autónoma de Sinaloa}\par
{\bfseries Facultad de Contaduría y Administración}\par
{\bfseries Licenciatura en Contaduría Pública — Modalidad Virtual}\par
\vspace{.5cm}
{\Large\bfseries Producto integrador}\par
{\large Ponencia sobre algún tema de las áreas administrativas}\par
\vspace{.5cm}
Martín Jonathan de la Cruz Muñoz\par
Nadia Aileen Valdez Acosta\par
GRUPO 103 LCP\par
13 de septiembre de 2026
\end{center}

\section{Introducción}
Completar tema, objetivos, tesis y fuentes principales.

\section{Desarrollo}
\subsection{Antecedentes conceptuales}
Completar.

\subsection{Precisión teórica}
Completar.

\subsection{Análisis y discusión}
Completar.

\section{Conclusiones}
Las conclusiones deben introducirse con un párrafo breve y después enumerarse en orden de importancia.

\section{Referencias}
% Completar con formato APA.

\bibliographystyle{plainnat}
\bibliography{referencias-investigacion-aplicada-a-la-contaduria/actividad-final/reporte-actividad-final-materialidad-operaciones-facturadas}
\end{document}